\subsection{Normal-Derived Distributions}
% Source integration: supplied Chapter 6 screenshots, Sections 6.1--6.4.

Independent $N(0,1)$ variables generate the normal-derived inference laws used below: $\chi^2$ sums, Student $t$ ratios, and $F$ ratios.

\begin{proposition} \label{prop:chi-square-normal-sums} \hlt{Chi-Square Sums}\\
Let $n\geq1$ and $Z_1,\ldots,Z_n\overset{\mathrm{ind}}{\sim}N(0,1)$.
\begin{enumerate}[label=(\roman*)]
\setlength{\itemsep}{0pt}
\item For $V=\sum_{i=1}^nZ_i^2$, $V\sim\chi_n^2=\operatorname{Gamma}(n/2,1/2)$ in the rate convention,\\
$f_V(v)=2^{-n/2}\Gamma(n/2)^{-1}v^{n/2-1}e^{-v/2}\mathbf{1}_{(0,\infty)}(v)$,\\
$\psi_V(t)=(1-2t)^{-n/2}$ for $t<1/2$, $\mathbb{E}V=n$, and $\operatorname{Var}(V)=2n$.
\item If $U\sim\chi_m^2$, $W\sim\chi_n^2$, and $U\perp\!\!\!\perp W$, then $U+W\sim\chi_{m+n}^2$.
\end{enumerate}
\begin{proof}
\begin{enumerate}[label=(\roman*)]
\setlength{\itemsep}{0pt}
\item
For $Z\sim N(0,1)$ and $t<1/2$, $\mathbb{E}e^{tZ^2}=(2\pi)^{-1/2}\int_{\mathbb{R}}e^{-(1-2t)z^2/2}\,dz=(1-2t)^{-1/2}$, so independence gives $\psi_V(t)=\prod_{i=1}^n\mathbb{E}e^{tZ_i^2}=(1-2t)^{-n/2}$.
MGF uniqueness gives $V\sim\operatorname{Gamma}(n/2,1/2)$, yielding the density; differentiating $\psi_V$ at $0$ gives $\mathbb{E}V=n$ and $\operatorname{Var}(V)=2n$.
\item If $U\perp\!\!\!\perp W$, then $\psi_{U+W}(t)=(1-2t)^{-m/2}(1-2t)^{-n/2}=(1-2t)^{-(m+n)/2}$, hence $U+W\sim\chi_{m+n}^2$.
\end{enumerate}
\hlt{Structural Commentary}: Chi-square degrees of freedom add by MGF exponents.\\
\hlt{Key Insight}: Squared independent standard normals form gamma sums.
\end{proof}
\end{proposition}

\begin{definition} \label{def:student-t-distribution} \hlt{Student $t$ Distribution}\\
For $n\geq1$, if $Z\sim N(0,1)$, $U\sim\chi_n^2$, and $Z\perp\!\!\!\perp U$, then $T=Z/\sqrt{U/n}$ has the \hlt{Student $t$ distribution} with $n$ degrees of freedom, written $T\sim t_n$.
\end{definition}

\begin{proposition} \label{prop:student-t-density} \hlt{Student $t$ Density}\\
Let $n\geq1$ and $T=Z/\sqrt{U/n}$ with $Z\sim N(0,1)$, $U\sim\chi_n^2$, and $Z\perp\!\!\!\perp U$.
\begin{enumerate}[label=(\roman*)]
\setlength{\itemsep}{0pt}
\item For $t\in\mathbb{R}$, $f_T(t)=\frac{\Gamma((n+1)/2)}{\sqrt{n\pi}\Gamma(n/2)}(1+t^2/n)^{-(n+1)/2}$.
\item For all $t\in\mathbb{R}$, $f_T(-t)=f_T(t)$.
\item For each fixed $t\in\mathbb{R}$, $\lim_{n\to\infty}f_T(t)=\phi(t)=(2\pi)^{-1/2}e^{-t^2/2}$.
\end{enumerate}
\begin{proof}
\begin{enumerate}[label=(\roman*)]
\setlength{\itemsep}{0pt}
\item For fixed $u>0$, $z=t\sqrt{u/n}$ has Jacobian $\sqrt{u/n}$, so $f_T(t)=\int_0^\infty\phi(t\sqrt{u/n})\sqrt{u/n}f_U(u)\,du$.\\
Substituting $\phi(z)=(2\pi)^{-1/2}e^{-z^2/2}$ and $f_U(u)=2^{-n/2}\Gamma(n/2)^{-1}u^{n/2-1}e^{-u/2}\mathbf{1}_{(0,\infty)}(u)$ gives\\
$f_T(t)=\frac{1}{\sqrt{n\pi}2^{(n+1)/2}\Gamma(n/2)}\int_0^\infty u^{(n+1)/2-1}e^{-(1+t^2/n)u/2}\,du$.\\
The integral equals $\Gamma((n+1)/2)2^{(n+1)/2}(1+t^2/n)^{-(n+1)/2}$, yielding the density.
\item By (i), $f_T(t)=c_nh_n(t^2)$ for $c_n>0$, so $f_T(-t)=f_T(t)$ for all $t$.
\item For fixed $t$, $f_T(t)=(2\pi)^{-1/2}\frac{\Gamma((n+1)/2)}{\sqrt{n/2}\Gamma(n/2)}(1+t^2/n)^{-(n+1)/2}$,\\
where the gamma ratio tends to $1$ and the last factor tends to $e^{-t^2/2}$.
\end{enumerate}
\hlt{Structural Commentary}: Student $t$ is a normal scale mixture with random chi-square scale.\\
\hlt{Key Insight}: Random denominator uncertainty thickens the normal tails.
\end{proof}
\end{proposition}

\begin{example} \label{ex:t-one-is-cauchy} \hlt{$t_1$ is Cauchy}\\
If $T\sim t_1$, then $f_T(t)=\Gamma(1)(\sqrt{\pi}\Gamma(1/2))^{-1}(1+t^2)^{-1}=1/(\pi(1+t^2))$, so $T\sim\operatorname{Cauchy}(0,1)$.
\end{example}

\begin{definition} \label{def:f-distribution} \hlt{$F$ Distribution}\\
For $m,n\geq1$, if $U\sim\chi_m^2$, $V\sim\chi_n^2$, and $U\perp\!\!\!\perp V$, then $W=(U/m)/(V/n)$ has the \hlt{$F$ distribution} with $(m,n)$ degrees of freedom, written $W\sim F_{m,n}$.
\end{definition}

\begin{proposition} \label{prop:f-density} \hlt{$F$ Density}\\
Let $m,n\geq1$ and $W=(U/m)/(V/n)$, where $U\sim\chi_m^2$, $V\sim\chi_n^2$, and $U\perp\!\!\!\perp V$.
\begin{enumerate}[label=(\roman*)]
\setlength{\itemsep}{0pt}
\item
$f_W(w)=\frac{\Gamma((m+n)/2)}{\Gamma(m/2)\Gamma(n/2)}(m/n)^{m/2}w^{m/2-1}(1+mw/n)^{-(m+n)/2}\mathbf{1}_{(0,\infty)}(w)$.
\item
If $n>2$, then $\mathbb{E}W=n/(n-2)$.
\end{enumerate}
\begin{proof}
\begin{enumerate}[label=(\roman*)]
\setlength{\itemsep}{0pt}
\item For $w>0$, set $u=mwv/n$; the Jacobian is $mv/n$, so $f_W(w)=\int_0^\infty f_U(mwv/n)f_V(v)(mv/n)\,dv$.\\
Substitution gives $f_W(w)=Cw^{m/2-1}\int_0^\infty v^{(m+n)/2-1}e^{-(1+mw/n)v/2}\,dv$,\\
where $C=(m/n)^{m/2}/(2^{(m+n)/2}\Gamma(m/2)\Gamma(n/2))$, and the gamma integral gives (i).
\item
Since $U\perp\!\!\!\perp V$, $\mathbb{E}W=(n/m)\mathbb{E}U\,\mathbb{E}(V^{-1})=n\,\mathbb{E}(V^{-1})=n/(n-2)$ for $n>2$.
\end{enumerate}
\hlt{Structural Commentary}: An $F$ variable is a ratio of independent chi-square averages.\\
\hlt{Key Insight}: Denominator variability controls the right tail and the mean condition $n>2$.
\end{proof}
\end{proposition}

\begin{corollary} \label{cor:f-reciprocal-and-squared-t} \hlt{Reciprocal $F$ and Squared $t$}
\begin{enumerate}[label=(\roman*)]
\setlength{\itemsep}{0pt}
\item
For $m,n\geq1$, if $W\sim F_{m,n}$, then $W^{-1}\sim F_{n,m}$.
\item
If $T\sim t_n$, then $T^2\sim F_{1,n}$.
\end{enumerate}
\begin{proof}
\begin{enumerate}[label=(\roman*)]
\setlength{\itemsep}{0pt}
\item
Write $W=(U/m)/(V/n)$ with $U\sim\chi_m^2$, $V\sim\chi_n^2$, and $U\perp\!\!\!\perp V$; then $W^{-1}=(V/n)/(U/m)\sim F_{n,m}$.
\item
Write $T=Z/\sqrt{V/n}$ with $Z\sim N(0,1)$, $V\sim\chi_n^2$, and $Z\perp\!\!\!\perp V$; since $Z^2\sim\chi_1^2$, $T^2=(Z^2/1)/(V/n)\sim F_{1,n}$.
\end{enumerate}
\hlt{Structural Commentary}: Reciprocals swap $F$ degrees of freedom, and squaring $t$ removes the normal sign.\\
\hlt{Key Insight}: Both identities are immediate from the same chi-square-ratio representation.
\end{proof}
\end{corollary}
